%!TEX program = pdflatex
% Mathematics article -- amsart, the neutral class of the subject. Drafting
% here and converting on acceptance costs nothing, because a journal class
% changes the front matter and the page layout and leaves the body alone.
\documentclass[11pt,reqno]{amsart}

% amsart's own page is set for the AMS printed page and reads narrow on
% screen; a journal class resets it at submission.
\usepackage[margin=1in]{geometry}

% amsart already loads amsmath, amsthm and amsfonts.
\usepackage{amssymb,mathtools}
\usepackage{graphicx}
\usepackage{bm}
\usepackage{booktabs}
\usepackage{xcolor}
\usepackage{hyperref}
% Links are coloured rather than boxed: the default frames print as
% rectangles around every citation and cross-reference.
\hypersetup{colorlinks=true,
            linkcolor=blue!55!black,
            citecolor=blue!55!black,
            urlcolor=blue!55!black}
\usepackage[capitalize,nameinlink]{cleveref}

% One shared counter, so a Lemma 2.3 and a Theorem 2.4 never collide.
\theoremstyle{plain}
\newtheorem{theorem}{Theorem}[section]
\newtheorem{lemma}[theorem]{Lemma}
\newtheorem{proposition}[theorem]{Proposition}
\newtheorem{corollary}[theorem]{Corollary}
\newtheorem{conjecture}[theorem]{Conjecture}

\theoremstyle{definition}
\newtheorem{definition}[theorem]{Definition}
\newtheorem{example}[theorem]{Example}

\theoremstyle{remark}
\newtheorem{remark}[theorem]{Remark}

% amsart numbers equations straight through; section numbering matches the
% theorem counter and survives inserting a section.
\numberwithin{equation}{section}

\title{Uniqueness of the Renormalized Quartic-Oscillator Critical Point}

\author{Zhiyuan Yao}
\address{Lanzhou Center for Theoretical Physics, Key Laboratory of Theoretical
  Physics of Gansu Province, Key Laboratory of Quantum Theory and Applications
  of MoE, Gansu Provincial Research Center for Basic Disciplines of Quantum
  Physics, Lanzhou University, Lanzhou 730000, China}
\email{yaozy@lzu.edu.cn}

\date{\today}

% AMS journals require a 2020 Mathematics Subject Classification, primary
% first; keywords are optional and printed with it.
% \subjclass[2020]{Primary 00A00; Secondary 00B00}
% \keywords{KEYWORDS}

\begin{document}

% amsart typesets the abstract from the title block, so it is
% declared inside the document and before that block.
\begin{abstract}
For the ground state $G_b$ of
$h_b=\tfrac12(-d^2/dx^2+bx^2)+x^4$, define
$m(b)=b+12\langle G_b,x^2G_b\rangle$.  We prove that $m'$ has exactly one
real zero $b_*$.  Moreover, $-16/5<b_*<0$ and $m''(b_*)>0$.  Exact dilation
reduces the problem to the level equation
$\lambda_1''(\alpha)=4/3$ for the first Montgomery eigenvalue.  Analytic
estimates confine every target to $0<\alpha<4/5$.  At any target, a reduced
resolvent identity turns an even-sector gap bound into strict downward
crossing.  The required gap bound is certified by exact moment recurrences
and an exhaustive rational Bernstein subdivision of a three-dimensional
box.  A translated Gaussian forces existence, and a zero-topology lemma
then gives uniqueness.  All load-bearing computer-assisted inequalities use
exact rational arithmetic.
\end{abstract}

\maketitle

\section{Introduction}

The renormalized oscillator coefficient $m$ was introduced by Baez and Zhou,
who conjectured from finite-difference evidence that it has one critical
point \cite{Baez_1992ro}.  Its quartic spectral background includes the
classical perturbative and analytic work of Bender--Wu and Simon
\cite{Bender_1969ao,Simon_1970ca}.  Analytic perturbation theory, the
standard theory of discrete Schr\"odinger spectra, and exponential decay of
bound states provide the operator-theoretic setting used below
\cite{Kato_Book,ReedSimonIV,Agmon_Book}.

Under the exact dilation used below, the spectral problem becomes the
Montgomery family
\[
H_\alpha=-\frac{d^2}{dt^2}+\left(\frac{t^2}{2}-\alpha\right)^2,
\]
which has a substantial magnetic-spectral history beginning with Montgomery
and continuing through magnetic-well and surface-superconductivity analysis
\cite{Montgomery_1995ht,Helffer_1996sa,Pan_2002so,HELFFER_2004mb,Fournais_Book}.
Helffer revisited the $k=1$ model, and Helffer--Persson treated the associated
higher-order anharmonic family \cite{Helffer_2010tm,Helffer_2010sp}.
Recent work connects quartic spectral summability to the Engel group
\cite{Bahouri_2023ss}, and the anharmonic oscillator itself remains a live
subject \cite{Turbiner_2021ao}.  Helffer and Léautaud proved uniqueness of the zero
of $\lambda_1'(\alpha)$ \cite{Helffer_2024oc}; that theorem does not decide
the present problem, whose exact stationarity condition is the
second-derivative level equation $\lambda_1''(\alpha)=4/3$.

Our proof separates global analysis from a finite exact certificate.  First,
all solutions of the level equation are placed in a compact rational
interval.  Second, reduced-resolvent identities show that strict
transversality follows from an even-gap inequality.  Moment and hypervirial
constraints reduce failure of that inequality to the simultaneous
nonnegativity of nine rational polynomials on a compact box.  The Bernstein
basis and its subdivision properties are surveyed in
\cite{Farouki_2012tb}; interval and verification methods provide the wider
validated-computation context \cite{Moore_Book,Rump_2010vm}.  Here every
Bernstein coefficient and every de Casteljau subdivision is instead computed
in exact rational arithmetic.  The resulting exhaustive tree excludes the
box in 757 nodes, so the certificate is part of the proof rather than a
floating-point experiment.

\begin{theorem}\label{thm:unique}
Let $h_b$ be the standard form-sum quartic oscillator on $L^2(\mathbb R)$,
let $G_b$ be its normalized positive ground state, and define
\[
m(b)=b+12\langle G_b,x^2G_b\rangle.
\]
Then $m$ has exactly one real critical point.  It lies in
$(-16/5,0)$ and is a strict local minimum.
\end{theorem}

\begin{remark}
Only the rational Bernstein exclusion in the transversality argument is
computer assisted.  The target localization, existence argument, and
zero-counting step are analytic.  All nine polynomials, their affine cube
map, the subdivision rule, the strict terminal test, and the exact coverage
argument are stated below.
\end{remark}

\section{Reduction to a curvature level set}
\label{sec:reduction}

\begin{proposition}
  \label{prop:reduction}
  Let $E(b)$ be the ground eigenvalue of $h_b$ and let
  \begin{equation}
    \label{eq:Halpha}
    H_\alpha=-\frac{d^2}{dt^2}+\left(\frac{t^2}{2}-\alpha\right)^2,
    \qquad
    \lambda(\alpha)=\inf\operatorname{spec}H_\alpha .
  \end{equation}
  Then $\lambda$ is real-analytic on $\mathbb R$, and with
  \begin{equation}
    \label{eq:f}
    f(\alpha):=\lambda''(\alpha)-\frac43
  \end{equation}
  one has
  \begin{equation}
    \label{eq:mprime}
    m'(b)=\frac32\,f(-b/4).
  \end{equation}
  Consequently $m$ has exactly one critical point if and only if $f$ has
  exactly one real zero.
\end{proposition}

\begin{proof}
  The Feynman--Hellmann formula applied to $\partial_bh_b=x^2/2$ gives
  \begin{equation}
    \label{eq:fh}
    m(b)=b+24E'(b),
    \qquad
    m'(b)=1+24E''(b).
  \end{equation}
  The unitary dilation $t=\sqrt2\,x$ together with $\alpha=-b/4$ carries
  $h_b$ onto $H_\alpha-\alpha^2$, so $E(b)=\lambda(-b/4)-b^2/16$.

  For the analyticity, write
  $H_\alpha=p^2+t^4/4-\alpha t^2+\alpha^2$. Multiplication by $t^2$ is
  infinitesimally form-bounded with respect to $p^2+t^4/4$, so the forms are
  closed on the common domain
  \begin{equation*}
    \mathcal Q=H^1(\mathbb R)\cap L^2(\mathbb R,t^4\,dt),
  \end{equation*}
  and $\{H_\alpha\}$ is an analytic family of type (B) with compact
  resolvent. In one dimension the ground eigenvalue is simple by
  Sturm--Liouville theory, so the local analytic perturbation branches
  concatenate into a real-analytic $\lambda$ on all of $\mathbb R$; in
  particular $f$ is $C^1$.

  Substituting $E(b)=\lambda(-b/4)-b^2/16$ into \eqref{eq:fh},
  \begin{equation*}
    m'(b)=-2+\frac32\lambda''(-b/4)=\frac32 f(-b/4),
  \end{equation*}
  which is \eqref{eq:mprime}. Since $b\mapsto-b/4$ is a bijection of
  $\mathbb R$, the zeros of $m'$ and of $f$ correspond one to one.
\end{proof}

Call $\alpha$ a \emph{target} when $f(\alpha)=0$. The rest of the paper
shows that there is exactly one target, and that it lies in $(0,4/5)$.

\section{Every target lies in \texorpdfstring{$(0,4/5)$}{(0,4/5)}}
\label{sec:localization}

\subsection{The nonpositive half-line}

\begin{proposition}
  \label{prop:nonpositive}
  $f(\alpha)>0$ for every $\alpha\le0$. In particular no target is
  nonpositive.
\end{proposition}

\begin{proof}
  Put $\alpha=-s$ with $s\ge0$, remove the additive constant $s^2$, and set
  \begin{equation*}
    K_s=p^2+s\,t^2+\frac{t^4}{4},
    \qquad
    e_j(s)=\lambda_j(-s)-s^2 .
  \end{equation*}
  Let $u$ be the normalized ground state of $K_s$ and put
  \begin{equation*}
    a=\langle t^2\rangle,\quad
    c=\langle t^4\rangle,\quad
    v=c-a^2,\quad
    g=(t^2-a)u,
  \end{equation*}
  with $R$ the positive reduced resolvent above $e_1$. Differentiating the
  eigenvalue equation twice gives
  \begin{equation}
    \label{eq:L1}
    \lambda''(-s)=2-2S,
    \qquad S=\langle g,Rg\rangle .
  \end{equation}
  Only even excited states contribute to $S$, so with $\Delta=e_3-e_1$ the
  spectral theorem gives $S\le v/\Delta$. The virial identity
  $2\langle p^2\rangle=2sa+c$ yields
  \begin{equation}
    \label{eq:L2}
    e_1=2sa+\frac34c\ \ge\ \frac34c,
    \qquad
    v\le c\le\frac43e_1 .
  \end{equation}

  It remains to separate $e_3$ from $e_1$. For $d>0$ let $r_d(s)>0$ solve
  \begin{equation}
    \label{eq:L3}
    s=r_d^2-\frac{d}{4r_d},
    \qquad
    F_d(s)=r_d-\frac{d}{16r_d^2}.
  \end{equation}
  A centred Gaussian of width $\rho$ has Rayleigh quotient
  $\rho/2+s/(2\rho)+3/(16\rho^2)$, whose stationarity equation is
  \eqref{eq:L3} with $d=3$; hence $e_1\le F_3$. For every $\gamma>s$,
  \begin{equation*}
    \frac{t^4}{4}+s\,t^2\ \ge\ \gamma t^2-(\gamma-s)^2 ,
  \end{equation*}
  and comparison with $p^2+\gamma t^2$, optimized at
  $\sqrt\gamma=r_5(s)$, gives $e_3\ge5F_5$. At fixed $s$, implicit
  differentiation of \eqref{eq:L3} gives
  \begin{equation}
    \label{eq:L4}
    \frac{\partial F_d(s)}{\partial d}=\frac1{16\,r_d(s)^2}>0 ,
  \end{equation}
  so $e_3\ge5F_5>5F_3\ge5e_1$ and therefore $\Delta>4e_1$. With
  \eqref{eq:L1} and \eqref{eq:L2},
  \begin{equation}
    \label{eq:S-third}
    S<\frac{(4/3)e_1}{4e_1}=\frac13,
    \qquad\text{hence}\qquad
    f(\alpha)=\lambda''(\alpha)-\frac43>0 . \qedhere
  \end{equation}
\end{proof}

\subsection{Necessary inequalities at a target}

\begin{lemma}
  \label{lem:target-ineq}
  Let $\alpha$ be a target, let $u$ be the corresponding normalized ground
  state, and write $a=\langle t^2\rangle$, $v=\langle t^4\rangle-a^2$. Then
  \begin{align}
    \label{eq:L5}
    v^2&\le\frac43a,
    &&\text{so } v\le U(a):=2\sqrt{\frac a3},\\
    \label{eq:pos}
    3a-2\alpha&>0,\\
    \label{eq:L6}
    v&\ge L(a,\alpha):=2\alpha a-a^2+\frac1{2a},\\
    \label{eq:L7}
    P(a,\alpha,v)&:=
      v\left[3a^3-6\alpha a^2+(3a+2\alpha)v+\frac92\right]-10a^2\ \ge\ 0 .
  \end{align}
\end{lemma}

\begin{proof}
  Inequality \eqref{eq:L5} is the inverse- and energy-weighted response
  identity at a target; \eqref{eq:pos} is the momentum double commutator;
  \eqref{eq:L6} combines the virial identity with the Heisenberg
  inequality; and \eqref{eq:L7} is Robertson's inequality for the pair
  $t^2$ and $D=(tp+pt)/2$, followed by the $t^2$ hypervirial identity.
\end{proof}

\begin{lemma}
  \label{lem:lambda-prime}
  $\lambda'(\alpha)=2\alpha-a>0$ for every $\alpha\ge1$; in particular
  $\alpha>a/2$ there.
\end{lemma}

\begin{proof}
  At a positive critical point of $\lambda$, half-line integration by parts
  gives
  \begin{equation}
    \label{eq:L8}
    (\lambda-\alpha^2)\,u(0)^2
    =\int_0^\infty (t+\sqrt{2\alpha})(t-\sqrt{2\alpha})^2u(t)^2\,dt>0 .
  \end{equation}
  A centred Gaussian of width $\rho=(2\alpha)^{-1}$ gives
  $\lambda(\alpha)\le1/(4\alpha)+\tfrac34\alpha^2\le\alpha^2$ for
  $\alpha\ge1$, contradicting \eqref{eq:L8}. So $\lambda'$ has no zero on
  $[1,\infty)$, and the well-bracketing asymptotic
  $\lambda(\alpha)\sim\sqrt{2\alpha}$ selects the positive sign. The
  half-line identity \eqref{eq:L8} is due to Pan and Kwek \cite{Pan_2002so}
  for the first eigenvalue, and the proof recalled here is that of Helffer and
  Persson Sundqvist \cite{Helffer_2010sp}; the comparison argument, the sign
  selection and the asymptotic are taken from Helffer and L\'eautaud
  \cite{Helffer_2024oc} and are not derived here.
\end{proof}

\begin{proposition}
  \label{prop:below-one}
  Every target satisfies $\alpha<1$.
\end{proposition}

\begin{proof}
  Suppose a target had $\alpha\ge1$. Then \eqref{eq:pos} and
  \Cref{lem:lambda-prime} give $a>2/3$, so $U(a)<3a^2$ and $P$ is
  decreasing in $\alpha$. Using \eqref{eq:L5}, \eqref{eq:L7} and
  \Cref{lem:lambda-prime},
  \begin{equation*}
    10a^2<U(a)\left(4aU(a)+\frac92\right)
    =\frac{16}{3}a^2+9\sqrt{\frac a3},
  \end{equation*}
  whence
  \begin{equation}
    \label{eq:L10}
    a<a_0:=\left(\frac{9\sqrt3}{14}\right)^{2/3}<\frac98 .
  \end{equation}
  On the other hand \eqref{eq:L5} and \eqref{eq:L6} give
  \begin{equation}
    \label{eq:L11}
    \alpha\le F(a):=\frac a2+\frac1{\sqrt{3a}}-\frac1{4a^2}.
  \end{equation}
  Writing $y=a^{3/2}$,
  \begin{equation*}
    2F'(a)=\frac{\bigl(y-1/(2\sqrt3)\bigr)^2+11/12}{y^2}>0,
  \end{equation*}
  so $F$ is increasing, and
  \begin{equation*}
    F(9/8)=\frac{473}{1296}+\sqrt{\frac8{27}}<\frac25+\frac35=1 .
  \end{equation*}
  Together with \eqref{eq:L10} this contradicts $\alpha\ge1$.
\end{proof}

\subsection{Sharpening the endpoint to \texorpdfstring{$4/5$}{4/5}}

\begin{lemma}
  \label{lem:L12}
  At a target, a unit-width Gaussian and the virial identity give
  \begin{equation}
    \label{eq:L12}
    v\le E(a,\alpha):=-a^2+\frac83\alpha a-\frac23\alpha+\frac{11}{12},
  \end{equation}
  and for $0<\alpha<1$ the polynomial $P$ of \eqref{eq:L7} is strictly
  increasing in $v>0$, because
  \begin{equation}
    \label{eq:L13}
    3a^3-6\alpha a^2+\frac92>\frac{17}{18}+\frac{(3a-4)^2(3a+2)}9>0 .
  \end{equation}
\end{lemma}

\begin{proposition}
  \label{prop:below-45}
  Every target satisfies $0<\alpha<4/5$.
\end{proposition}

\begin{proof}
  By \Cref{prop:nonpositive,prop:below-one} a target lies in $(0,1)$;
  suppose one had $\alpha\ge4/5$. From $L\le U$ one gets $\alpha\le F(a)$
  with $F$ as in \eqref{eq:L11}, and $F$ is increasing with
  \begin{equation*}
    F(19/20)=\frac{19}{40}+\sqrt{\frac{20}{57}}-\frac{100}{361}<\frac45,
  \end{equation*}
  the radical comparison following exactly from
  $(8693/14440)^2-20/57=7216747/625540800>0$. Hence $a>19/20$, so
  $U(a)<3a^2$ and $\partial_\alpha P=2v(v-3a^2)<0$.

  \emph{Case $19/20<a\le2$.} Put $x=\sqrt{a/3}$, so $9/16<x<5/6$.
  Monotonicity in $v$ (\Cref{lem:L12}) and in $\alpha$ gives
  $P(a,\alpha,v)\le P(a,4/5,2x)=\tfrac x5Q(x)$ with
  \begin{equation}
    \label{eq:Q}
    Q(x)=810x^6-432x^4-270x^3+32x+45 .
  \end{equation}
  The Sturm chain of $Q$ has three sign variations at both $9/16$ and
  $5/6$, so $Q$ has no zero between them, and $Q(3/4)=-76659/2048<0$. Thus
  $Q<0$ on the whole range, contradicting $P\ge0$.

  \emph{Case $a\ge2$.} Positivity of $v$ together with \eqref{eq:L12} and
  $\alpha<1$ confines this case to $2\le a<3$. Since $P$ increases in $v$,
  \begin{equation*}
    P(a,\alpha,v)\le R(a,\alpha):=P\bigl(a,\alpha,E(a,\alpha)\bigr).
  \end{equation*}
  Exact differentiation gives
  \begin{equation}
    \label{eq:L15}
    \partial_\alpha^2R=\frac{4(4a-1)}9\,J(a,\alpha),
    \qquad
    J=-6a^2+48a\alpha-6a-12\alpha+11 .
  \end{equation}
  On $2\le a\le3$, $4/5\le\alpha\le1$ one has $J>0$: it increases in
  $\alpha$, and the concave quadratic $J(a,4/5)$ has endpoint values
  $211/5$ and $223/5$. Also
  \begin{equation*}
    \partial_\alpha R(a,4/5)=\frac{S(a)}{1800},
    \qquad
    S=-11760a^3+47532a^2+22424a-6343,
  \end{equation*}
  and $S$ is concave on $[2,3]$ with positive endpoint values $134553$ and
  $171197$, so $R$ increases in $\alpha$. Finally
  \begin{equation*}
    R(a,1)=-\frac{K(a)}{144},
    \qquad
    K=876a^3-176a^2-2139a-180,
  \end{equation*}
  with $K(2)=1846$, $K'(2)=7669$ and $K''(a)=5256a-352>0$ for $a\ge2$, so
  $K>0$ and $R(a,\alpha)<R(a,1)<0$, again contradicting $P\ge0$.
\end{proof}

\section{Strict transversality at every target}
\label{sec:transversality}

\subsection{A sufficient even-gap criterion}

\begin{lemma}
  \label{lem:gap-criterion}
  Let $\alpha$ be a target with normalized ground state $u$, and set
  $a=\langle t^2\rangle$, $g=(t^2-a)u$, $R=Q(H_\alpha-\lambda_1)^{-1}Q$,
  $\phi=Rg$, and $\Delta=\lambda_3-\lambda_1$ for the first even gap. Then
  \begin{equation}
    \label{eq:T1}
    \langle g,Rg\rangle=\frac13,
    \qquad
    \lambda'''(\alpha)=-6\langle\phi,(t^2-a)\phi\rangle,
  \end{equation}
  and
  \begin{equation}
    \label{eq:transversality}
    \Delta>3a^2
    \quad\Longrightarrow\quad
    \lambda'''(\alpha)<0 .
  \end{equation}
\end{lemma}

\begin{proof}
  Differentiating in the parallel-transport gauge gives \eqref{eq:T1}; the
  first identity is the target condition $\lambda''=4/3$ rewritten through
  \eqref{eq:L1}. The minimax principle and the Cauchy--Schwarz inequality
  give respectively
  \begin{equation}
    \label{eq:T2}
    \lVert\phi\rVert^2\le\frac1{3\Delta},
    \qquad
    \langle\phi,t^2\phi\rangle\ge\frac1{9a},
  \end{equation}
  so that
  $\langle\phi,(t^2-a)\phi\rangle\ge(\Delta-3a^2)/(9a\Delta)$, which is
  positive exactly when $\Delta>3a^2$.
\end{proof}

\subsection{Isolating the third pure-quartic eigenvalue}

\begin{lemma}
  \label{lem:quartic}
  Let $L=p^2+t^4/4$. Then $\lambda_3(L)>467/100$.
\end{lemma}

\begin{proof}
  Take the even trial vector $\psi(t)=P(t^2)e^{-t^2/2}$ with the following
  decimal strings read as exact rationals:
  \begin{equation}
    \label{eq:T3}
    \begin{aligned}
      P(y)={}&1-1.8476436929042896\,y-0.13114586665910161\,y^2
        +0.050796559364988544\,y^3\\
      &+0.0027570484421582321\,y^4-0.0012003885725975921\,y^5\\
      &+0.000098882369961766198\,y^6-0.0000033908436935998999\,y^7\\
      &+0.000000042707174142248135\,y^8 .
    \end{aligned}
  \end{equation}
  With $\mu=\langle\psi,L\psi\rangle/\lVert\psi\rVert^2$ and
  $r^2=\lVert(L-\mu)\psi\rVert^2/\lVert\psi\rVert^2$, Gaussian moments
  reduce both to exact rationals, and direct reduction verifies
  \begin{equation}
    \label{eq:T4}
    \frac{467}{100}<\mu<\frac{88439}{10000},
    \qquad
    \left(\mu-\frac{467}{100}\right)^2>r^2,
    \qquad
    \left(\frac{88439}{10000}-\mu\right)^2>r^2 .
  \end{equation}
  A unit Gaussian gives $\lambda_1(L)\le11/16$, and the comparison
  $t^4/4\ge(169/100)t^2-(169/100)^2$ gives
  $\lambda_5(L)\ge88439/10000$. By parity, the residual theorem applied to
  \eqref{eq:T4} isolates a single even eigenvalue in the stated interval,
  which is therefore $\lambda_3(L)$.
\end{proof}

\begin{corollary}
  \label{cor:G}
  At every target,
  \begin{equation}
    \label{eq:GG}
    \Delta-3a^2>G(\alpha,a,v):=
      \frac{467}{125}-\frac{125}{61}\alpha^2+2\alpha a
      -\frac{15}{4}a^2-\frac34v .
  \end{equation}
\end{corollary}

\begin{proof}
  For $\theta=(4/5)^3=64/125$, completing the square gives
  $t^4/4-\alpha t^2\ge\theta t^4/4-\alpha^2/(1-\theta)$ with
  $1-\theta=61/125$, so
  $H_\alpha\ge p^2+\theta t^4/4-\tfrac{64}{61}\alpha^2$. Quartic scaling
  multiplies the eigenvalues of $p^2+\theta t^4/4$ by
  $\theta^{1/3}=4/5$, so by \Cref{lem:quartic}
  \begin{equation}
    \label{eq:lam3}
    \lambda_3(\alpha)>\frac45\cdot\frac{467}{100}-\frac{64}{61}\alpha^2
    =\frac{467}{125}-\frac{64}{61}\alpha^2 .
  \end{equation}
  At a target the virial identity gives
  $\lambda_1=\tfrac34(a^2+v)-2\alpha a+\alpha^2$. Subtracting it from
  \eqref{eq:lam3} and then subtracting $3a^2$ produces \eqref{eq:GG}: the
  coefficient of $\alpha^2$ moves from $64/61$ to
  $64/61+1=125/61$ because of the $+\alpha^2$ in $\lambda_1$, and the
  coefficient of $a^2$ from $3$ to $3+\tfrac34=\tfrac{15}4$.
\end{proof}

\subsection{The bad set and its exact constraints}

Let $m_j=\langle t^j\rangle$, so $m_0=1$, $m_2=a$ and $m_4=a^2+v$.
Eigenvalue and hypervirial integration by parts give, for $k\ge1$,
\begin{equation}
  \label{eq:T5}
  m_{2k+4}=\frac{4(2k+1)}{2k+3}(\lambda_1-\alpha^2)m_{2k}
    +\frac{8\alpha(k+1)}{2k+3}m_{2k+2}
    +\frac{2k(2k-1)(2k+1)}{2k+3}m_{2k-2},
\end{equation}
which determines $m_6,\dots,m_{20}$ as rational polynomials in
$(\alpha,a,v)$. Beyond the four inequalities of \Cref{lem:target-ineq} and
\eqref{eq:L12}, three further constraints are expectations of squares,
\begin{equation}
  \label{eq:T7}
  \bigl\langle q_i(t^2)^2\bigr\rangle
  =\sum_{j,k=0}^5q_{i,j}q_{i,k}\,m_{2(j+k)}\ \ge\ 0,
  \qquad i=1,2,3,
\end{equation}
with
\begin{equation}
  \label{eq:qvec}
  \begin{aligned}
    10^5q_1&=(-13381,-31687,100000,-48651,7725,-379),\\
    10^5q_2&=(-15103,-58143,100000,-32684,2975,-35),\\
    10^5q_3&=(-485,-70487,100000,-37526,4978,-207).
  \end{aligned}
\end{equation}

\begin{lemma}
  \label{lem:box}
  Every target lies in the closed rational box
  \begin{equation}
    \label{eq:box}
    0\le\alpha\le\frac45,
    \qquad
    0\le a\le\frac52,
    \qquad
    0\le v\le2 .
  \end{equation}
\end{lemma}

\begin{proof}
  The first bound is \Cref{prop:below-45}. If $a\ge1/4$, the left side of
  \eqref{eq:L12} before subtracting $v$ is maximized at $\alpha=4/5$, is
  decreasing at and beyond $a=5/2$, and equals $-8/15$ there; hence
  $a<5/2$. The response inequality \eqref{eq:L5} then gives
  $v^2<10/3<4$.
\end{proof}

A hypothetical target at which $G\le0$ therefore makes nine rational
polynomials simultaneously nonnegative on the box \eqref{eq:box}: the four
of \Cref{lem:target-ineq}, the Gaussian--virial constraint \eqref{eq:L12},
the three squares \eqref{eq:T7}, and $-G$. Call the set of such points the
\emph{bad set}.

\subsection{Finite rational Bernstein exclusion}

Map \eqref{eq:box} onto $[0,1]^3$ by $\alpha=\tfrac45x$, $a=\tfrac52y$,
$v=2z$. For a power polynomial
$p=\sum_{\boldsymbol j\le\boldsymbol d}A_{\boldsymbol j}
x^{j_1}y^{j_2}z^{j_3}$ the tensor Bernstein coefficients are
\begin{equation}
  \label{eq:T8}
  B_{\boldsymbol i}
  =\sum_{\boldsymbol j\le\boldsymbol i}A_{\boldsymbol j}
   \prod_{\ell=1}^3
   \frac{\binom{i_\ell}{j_\ell}}{\binom{d_\ell}{j_\ell}},
\end{equation}
and the range of $p$ over the box lies in the convex hull of the
$B_{\boldsymbol i}$; exact de Casteljau bisection supplies the coefficients
of each child box. The traversal always bisects a least-refined coordinate,
breaking ties in the order $\alpha,a,v$, and closes a box only when the
greatest Bernstein coefficient of one required polynomial is strictly
negative. Carried out in exact rational arithmetic it terminates, with the
closures distributed as follows.

\begin{center}
  \begin{tabular}{lr}
    \toprule
    Excluding polynomial & Leaves\\
    \midrule
    $3a-2\alpha$ & 0\\
    $4a/3-v^2$ & 78\\
    Heisenberg--virial & 9\\
    Gaussian--virial & 14\\
    Robertson--hypervirial & 18\\
    $\langle q_1^2\rangle$ & 33\\
    $\langle q_2^2\rangle$ & 72\\
    $\langle q_3^2\rangle$ & 66\\
    $-G$ & 89\\
    \bottomrule
  \end{tabular}
\end{center}

\begin{proposition}
  \label{prop:bernstein}
  The bad set is empty: $G>0$ at every target.
\end{proposition}

\begin{proof}
  The traversal produces a finite rooted binary tree with $379$ terminal
  nodes, $378$ internal nodes, and maximum total depth $24$. At each
  terminal node, direct substitution in the rational de Casteljau
  recurrence gives a strictly negative greatest Bernstein coefficient for
  the polynomial named in the table, so that polynomial is negative
  throughout the corresponding closed dyadic box and the box contains no
  point of the bad set.

  For exhaustiveness, label each child by appending $L$ or $R$ to its
  parent's binary path. The terminal paths are prefix-free, and their exact
  dyadic weights satisfy
  \begin{equation}
    \label{eq:kraft}
    \sum_{\text{terminal paths }w}2^{-|w|}=1 .
  \end{equation}
  Since every internal node has exactly the two de Casteljau children of
  its own box, prefix-freeness together with \eqref{eq:kraft} forces the
  terminal boxes to cover $[0,1]^3$, boundaries included. Every point of
  the root cube therefore lies in a terminal box, and every terminal box is
  excluded. The inputs to this finite calculation are \eqref{eq:T5},
  \eqref{eq:T7}, \eqref{eq:qvec}, \eqref{eq:T8}, the affine map above, and
  the stated deterministic split rule; no floating-point estimate and no
  stored polynomial enters.
\end{proof}

\begin{theorem}
  \label{thm:transversal}
  $f'(\alpha)=\lambda'''(\alpha)<0$ at every target.
\end{theorem}

\begin{proof}
  \Cref{prop:bernstein} gives $G>0$ at every target, \Cref{cor:G} then
  gives $\Delta>3a^2$, and \Cref{lem:gap-criterion} converts that into
  $\lambda'''(\alpha)<0$.
\end{proof}

\section{A target exists}
\label{sec:existence}

\begin{proposition}
  \label{prop:existence}
  $f$ has at least one positive zero.
\end{proposition}

\begin{proof}
  Suppose not. By \Cref{prop:nonpositive}, $f(0)>0$, so continuity and the
  absence of a zero give
  \begin{equation}
    \label{eq:no-zero}
    \lambda''(\alpha)>\frac43\qquad(\alpha\ge0),
  \end{equation}
  and integrating twice,
  \begin{equation}
    \label{eq:quadratic-growth}
    \lambda(\alpha)>\lambda(0)+\lambda'(0)\alpha+\frac23\alpha^2
    \qquad(\alpha>0).
  \end{equation}
  Against this stands a variational upper bound. For $\rho>0$ put
  $s=\sqrt{2\alpha}$ and take the normalized translated Gaussian
  \begin{equation*}
    \psi_{\rho,s}(t)=\left(\frac\rho\pi\right)^{1/4}
      \exp\left[-\frac\rho2(t-s)^2\right].
  \end{equation*}
  Writing $X=t-s$ one has $\langle X\rangle=\langle X^3\rangle=0$,
  $\langle X^2\rangle=1/(2\rho)$ and $\langle X^4\rangle=3/(4\rho^2)$, and
  $s^2=2\alpha$ gives $t^2/2-\alpha=sX+X^2/2$, so
  \begin{equation}
    \label{eq:gauss-moments}
    \langle p^2\rangle=\frac\rho2,
    \qquad
    \left\langle\left(\frac{t^2}{2}-\alpha\right)^2\right\rangle
    =\frac\alpha\rho+\frac3{16\rho^2}.
  \end{equation}
  Choosing $\rho=\sqrt{2\alpha}$ in the Rayleigh quotient,
  \begin{equation}
    \label{eq:upper}
    \lambda(\alpha)\le\sqrt{2\alpha}+\frac3{32\alpha}=o(\alpha^2),
  \end{equation}
  which contradicts \eqref{eq:quadratic-growth} as $\alpha\to\infty$.
\end{proof}

\section{The target is unique}
\label{sec:uniqueness}

\begin{lemma}
  \label{lem:zero-topology}
  Let $I\subset\mathbb R$ be an interval and $F\in C^1(I)$ with $F'(x)<0$ at
  every zero $x$ of $F$. Then $F$ has at most one zero in $I$.
\end{lemma}

\begin{proof}
  Suppose $x_1<x_2$ are zeros. The negative derivatives give $F<0$
  immediately to the right of $x_1$ and $F>0$ immediately to the left of
  $x_2$, so there are $p<q$ between them with $F(p)<0<F(q)$. Let $z$ be the
  first zero of $F$ in $[p,q]$. Continuity gives $F<0$ immediately to the
  left of $z$, whereas $F'(z)<0$ gives $F>0$ there --- a contradiction.
\end{proof}

\begin{proof}[Proof of \Cref{thm:unique}]
  By \Cref{prop:reduction} it suffices to show that $f$ has exactly one real
  zero. \Cref{thm:transversal} gives $f'<0$ at every zero of $f$, so
  \Cref{lem:zero-topology} leaves at most one; \Cref{prop:existence} gives at
  least one. Hence there is exactly one target $\alpha_*$, and
  \Cref{prop:nonpositive,prop:below-45} place it in $(0,4/5)$.

  Finally $b\mapsto\alpha=-b/4$ is a bijection and \eqref{eq:mprime} gives
  $m'(b)=0$ exactly when $f(-b/4)=0$, so $m$ has the single critical point
  $b_*=-4\alpha_*\in(-16/5,0)$, and
  \begin{equation*}
    m''(b_*)=-\frac38f'(\alpha_*)>0
  \end{equation*}
  by \Cref{thm:transversal}, making it a strict local minimum.
\end{proof}

\section{Scope and conclusion}

\Cref{thm:unique} concerns the ground eigenvalue and all real $b$.  It does
not locate $b_*$ sharply, address excited-state curvature levels, or settle
the infinite-mode field problem that motivated the oscillator.  The result
is logically distinct from uniqueness of a zero of $\lambda_1'$ and from
the separate positivity and coercivity problem for $m$.

\section*{Code availability}

One step of the proof is computer assisted: \Cref{prop:bernstein}, the
exclusion of the bad set, rests on a finite traversal that was carried out
rather than written out. An exact verifier that reproduces it in rational
arithmetic accompanies this paper; it rebuilds the nine polynomials from the
recurrence \eqref{eq:T5}, converts them by \eqref{eq:T8}, subdivides the box
\eqref{eq:box}, and prints the node, leaf, depth, and per-polynomial counts
reported in \Cref{sec:transversality}.

It is a check and not a premise. The nine polynomials, the affine map, the
deterministic split rule, the strict terminal test, and the prefix-free
coverage argument are all stated above, so the argument can be followed
without running anything, and no floating-point value enters at any point.

\section*{Acknowledgments}

Z.\ Y.\ acknowledges support from the Strategic Priority Research Program of the
Chinese Academy of Sciences under Grant No.~XDB1680102. Z.\ Y.\ also
acknowledges support by National Natural Science Foundation of China (Grant
No.~12247101). We also acknowledge Gewu Intelligence Lab for providing the
collaborative research environment in which this project was developed.

OpenAI GPT-5.6 Sol was used substantively in this work: in solving the
problem, including the formulation of arguments and the derivation of
intermediate steps, and in drafting and revising the manuscript. The
authors specified the problem, guided the approach to be taken, checked
every argument and result it produced, and take full responsibility for
the content of this manuscript.

\bibliographystyle{amsplain}
\bibliography{references,unverified}
\end{document}
